\section{The preparation counterexample}
\label{app:taylor}

For $u(0)=0$ the first Taylor coefficients of the backward flow are $c_1=d$,
$c_2=J_ed/2$, $c_3=(J_ec_2+.1D(d,d))/3$ and $c_4=(J_ec_3+.2D(d,c_2))/4$.
Substituting in the $RR$ coordinate of the two-site source gives
\[
 (c_4)_{RR}=-\frac{16087}{8000000}-\frac{29}{480000}\,e
            -\frac{29}{40000}\,e^2 .
\]
The coefficients through order three agree between $e=.01$ and $e=.3$, and
their fourth-order difference is the coefficient displayed in
\eqref{eq:taylor}:
\[
 (c_4)_{RR}\big|_{e=.3}-(c_4)_{RR}\big|_{e=.01}
 =-\frac{49619}{600000000}.
\]
Analyticity of a polynomial ODE near zero then proves the strict sign on some
positive interval, which is the content of \eqref{eq:taylor}. Numerically, at
$t=1$ the two extinction probabilities are $.2493000462$ and $.2492590475$, a
difference of about $4.1\cdot10^{-5}$; these displayed finite-time values are
diagnostics and not part of the analytic proof. The example shows that stronger
erasure can lower the extinction probability at short times from a
repressed preparation, because erasure also removes repressive marks and
permits access to the protected active states of this source. It is a
counterexample to universal monotonicity in the actuator, which is one reason
the amplitude certificate of Theorem~\ref{thm:amp} is stated against arbitrary
predictable policies rather than against constant actions.
